\documentclass[10pt,a4paper]{amsart}
\usepackage[margin=19mm]{geometry}
\usepackage{amsmath,amssymb,mathtools}
\usepackage[scr=rsfs,cal=euler]{mathalfa}
\usepackage{xfrac}
\usepackage[unicode,hidelinks,bookmarks=false]{hyperref}
\newcommand{\ie}{i.e.\ }
\newcommand{\cf}{cf.\ }
\newcommand{\resp}{respectively\ }
\newcommand{\Subsection}{\subsection}
\providecommand{\nobreakdash}{\nobreak\hbox{-}}
\setlength{\emergencystretch}{2em}
\multlinegap0pt
\usepackage{bm}
\usepackage{bbm}
\usepackage{dsfont}
\usepackage{xspace}
\usepackage{cancel}
\usepackage{mathtools}
\usepackage{amsthm}
\makeatletter
\@ifundefined{theorem}{\newtheorem{theorem}{Theorem}}{}
\makeatother
\title[Bonded braids and the Markov theorem]{Bonded braids and the Markov theorem}
\author{Paolo Cavicchioli, Bo\v{s}tjan Gabrov\v{s}ek, Matic Simoni\v{c}}
\date{Source: arXiv:2507.04565v2 (CC BY 4.0)}
\begin{document}
\maketitle

\noindent\emph{Source and licence.} Bonded braids and the Markov theorem. Version arXiv:2507.04565v2 (CC BY 4.0), distributed under the Creative Commons Attribution 4.0 International licence, as recorded in the arXiv licence metadata for this exact version and shown on the abstract page https://arxiv.org/abs/2507.04565v2. Full rights notice: licenses/arxiv-2507.04565-CC-BY-4.0.txt.\par\smallskip

\noindent\textit{Editorial scope.} Counted unit: the Theorem whose source label is \texttt{None} in the article (source0.tex lines 994--1058, source label \texttt{None}), delivered together with the article's own complete original proof of it (source0.tex lines 1060--1147, one \texttt{proof} environment of 88 lines). No mathematics is written, repaired or re-derived: statement and proof are the article's own text line for line. Required context retained uncounted: none. Every definition, display and label of those lines is retained unchanged. Omitted from this package and not redistributed: 1--993, 1059--1059, 1148--1549. Nothing inside a retained excerpt is omitted or altered: every retained body is delivered exactly as the article's lines are written, so no omission receipt of any kind accompanies this package. Citations: the retained excerpts carry no citation command at all, so no bibliography entry of the article is transcribed and no reference is redistributed. Reference adaptation: none was needed and none was made; every reference inside the delivered unit resolves inside it, so no unresolved reference or citation is produced. Printed slips kept literal: no printed word, symbol, hypothesis or number of the counted statement or of its proof was altered. Layout and class: the article's own class is \texttt{article} with packages including babel, inputenc, fontenc, geometry, amssymb, amsmath, amsthm, mathtools, epsfig, caption, todonotes, hyperref, xcolor, pinlabel; the wrapper is the lane's pinned \texttt{amsart} editorial wrapper at 10pt on A4, which restates the theorem-like environments the retained text needs and adds the article's own macro definitions that the retained text needs (none), with extra wrapper packages (bm, bbm, dsfont, xspace, cancel, mathtools). The delivered numbering is the unit's own. Assets: no retained span contains a figure environment, a graphics-inclusion command, a PDF-page inclusion, a table or a TikZ picture; no image file is copied into this package, the package declares no asset dependency and no image file is redistributed with it. Licence: version arXiv:2507.04565v2 of this article is distributed under the Creative Commons Attribution 4.0 International licence, as recorded in the arXiv licence metadata for this exact version and shown on the abstract page https://arxiv.org/abs/2507.04565v2. A full rights notice is delivered at licenses/arxiv-2507.04565-CC-BY-4.0.txt. Characters: every retained excerpt is pure ASCII, so the delivered engine, XeLaTeX with its default Latin Modern text font, reports no missing character.
\begin{theorem}
Let $n \geq 3$, and let $\sigma_i, b_i \in M_n$. The reduced bonded Burau representation is given explicitly by the mapping
for the braid generators:
\begin{align}
\sigma_1 &
\mapsto A_1' =
\begin{pmatrix}
-t & 0 & 0 \\
1 & 1 & 0 \\
0 & 0 & I_{n-3}
\end{pmatrix}, \qquad
\sigma_{n-1}
\mapsto A_{n-1}' =
\begin{pmatrix}
I_{n-3} & 0 & 0 \\
0 & 1 & t \\
0 & 0 & -t
\end{pmatrix},
\notag
\\
\sigma_{i} & \mapsto 
A_i' =
\begin{pmatrix}
I_{i-2} & 0 & 0 & 0 & 0 \\
0 & 1 & t & 0 & 0 \\
0 & 0 & -t & 0 & 0 \\
0 & 0 & 1 & 1 & 0 \\
0 & 0 & 0 & 0 & I_{n-i-2}
\end{pmatrix}\quad \text{for \; $1 < i < n-1$}. \notag
\end{align} 

and for the bond generators:

\begin{align}
b_1 &
\mapsto B_1' =
\begin{pmatrix}
-tz - z + 1 & 0 & 0 \\
z & 1 & 0 \\
0 & 0 & I_{n-3}
\end{pmatrix}, \qquad
b_{n-1}
\mapsto B_{n-1}' =
\begin{pmatrix}
I_{n-3} & 0 & 0 \\
0 & 1 & tz \\
0 & 0 & -tz - z + 1
\end{pmatrix},
\notag
\\
b_{i} &\mapsto 
B_i' =
\begin{pmatrix}
I_{i-2} & 0 & 0 & 0 & 0 \\
0 & 1 & tz & 0 & 0 \\
0 & 0 & -tz - z + 1 & 0 & 0 \\
0 & 0 & z & 1 & 0 \\
0 & 0 & 0 & 0 & I_{n-i-2}
\end{pmatrix}\quad \text{for \; $1 < i < n-1$}. \notag
\end{align} 


%and define matrices $A_1', A_2', \dots, A_{n-1}'$ and $B_1', B_2', \dots, B_{n-1}'$ in $\mathrm{GL}_{n-1}(\mathbb{Z}[t^{\pm 1}, z])$ as follows.

\end{theorem}

\begin{proof}
Let us first show how the reduced presentation is obtained. For all $i = 1, \dots, n-1$ we have:
$$
C^{-1} A_i C =
\begin{pmatrix}
A_i' & 0 \\
\ast_i & 1
\end{pmatrix},
$$
where $C = C_n$ is the $n \times n$ matrix
$$
C =
\begin{pmatrix}
1 & 1 & \cdots & 1 \\
0 & 1 & \cdots & 1 \\
0 & 0 & \ddots & \vdots \\
0 & 0 & \cdots & 1
\end{pmatrix},
$$
and $\ast_i$ is the row vector in $\mathbb{Z}^{n-1}$ given by:
$$
\ast_i =
\begin{cases}
(0, \dots, 0), & i < n-1 \\
(0, \dots, 0, 1), & i = n-1.
\end{cases}
$$
Similarly, for all $i = 1, \dots, n-1$ we have:
$$
C^{-1} B_i C =
\begin{pmatrix}
B_i' & 0 \\
\ast_i & 1
\end{pmatrix},
$$
with $\ast_i$ as above, but for $i = n-1$ given by $(0, \dots, 0, z)$. For each $i = 1, \dots, n-1$, define:
$$
A_i'' :=
\begin{pmatrix}
A_i' & 0 \\
\ast_i & 1
\end{pmatrix},
\quad
B_i'' :=
\begin{pmatrix}
B_i' & 0 \\
\ast_i & 1
\end{pmatrix}.
$$
In order to prove the theorem, it suffices to show that:
$$
A_i C = C A_i'', \quad \text{and} \quad B_i C = C B_i''.
$$
We observe that both $A_i C$ and $C A_i''$ are obtained from $C$ by the same row operations:
\begin{itemize}
    \item replace the entry in row $i$, column $i$, with $1 - t$,
    \item replace the entry in row $i+1$, column $i$, with $1$.
\end{itemize}
Similarly, $B_i C$ and $C B_i''$ are obtained from $C$ by:
\begin{itemize}
    \item replacing entry $(i,i)$ with $1 - tz$,
    \item replacing entry $(i+1,i)$ with $z$.
\end{itemize}

Thus, $A_i C = C A_i''$ and $B_i C = C B_i''$, proving the conjugation formulas.\\
\\
Since the matrices $A_1, \dots, A_{n-1}$ and $B_1, \dots, B_{n-1}$ satisfy the braid relations (with bonded extensions), their conjugates under $C^{-1} (-) C$ also satisfy these relations.

From the decomposition:
$$
C^{-1} A_i C = \begin{pmatrix} A_i' & 0 \\ \ast_i & 1 \end{pmatrix}, \quad
C^{-1} B_i C = \begin{pmatrix} B_i' & 0 \\ \ast_i & 1 \end{pmatrix},
$$
we see that the matrices $A_i'$ and $B_i'$ also satisfy braid-type relations.

Thus, we obtain reduced representation:
$$
\psi_n^r: M_n \longrightarrow \mathrm{GL}_{n-1}(\Lambda), \quad \sigma_i \mapsto A_i', \quad b_i \mapsto B_i',
$$
which we call the \emph{reduced bonded Burau representation}.

For $n = 2$, these reduce to:
$$
\psi_2^r : M_2 \longrightarrow \mathrm{GL}_1(\Lambda), \quad \sigma_1 \mapsto -t, \quad b_1 \mapsto -tz - z + 1.
$$


\end{proof}

\end{document}
